\section*{Hoofdstuk \ref{ch:b2:vertebrate-development} --- Embryonale ontwikkeling van gewervelden}

\begin{solution}{exo:b2:vertebrate-development:1}
\omterm{def:b2:vertebrate-development:egg}{Klieving}: snelle mitosen zonder groei, die de \omterm{def:b2:vertebrate-development:blastula}{blastula} opleveren.
\omterm{def:b2:vertebrate-development:blastula}{Blastula}: een holle bol of schijf van kleine cellen rond het
\omterm{def:b2:vertebrate-development:blastula}{blastocoel}. Gastrulatie: de celbewegingen die mesoderm en endoderm naar
binnen brengen en de drielagige gastrula met een darmholte en een chorda
opleveren. Neurulatie: het opvouwen van het dorsale ectoderm tot de
\omterm{prop:b2:vertebrate-development:neurulation}{neurale buis}, samen met de segmentatie van de somieten, wat de neurula met
het bouwplan van een gewervelde oplevert.
\end{solution}

\begin{solution}{exo:b2:vertebrate-development:2}
Ectoderm: opperhuid, hersenen en ruggenmerg, perifere zenuwen (\omterm{prop:b2:vertebrate-development:neurulation}{neurale
lijst}), lens en pigmentcellen. Mesoderm: chorda, spieren, bot en
kraakbeen, nieren, hart en bloed, lederhuid. Endoderm: de darmwand, lever,
alvleesklier, longen, schildklier.
\end{solution}

\begin{solution}{exo:b2:vertebrate-development:3}
Weinig \omterm{def:b2:vertebrate-development:egg}{dooier}: volledige, bijna gelijke \omterm{def:b2:vertebrate-development:egg}{klieving} en gastrulatie door
involutie via een ronde blastoporus in een holle bol (kikker: volledig
maar ongelijk, met grote en trage dooierrijke vegetatieve cellen). Veel
\omterm{def:b2:vertebrate-development:egg}{dooier} (kuiken): \omterm{def:b2:vertebrate-development:egg}{klieving} beperkt tot een schijf bovenop de \omterm{def:b2:vertebrate-development:egg}{dooier}, en
gastrulatie door een spleet, de primitiefstreep, in de platte schijf,
waarbij de lagen zich over de \omterm{def:b2:vertebrate-development:egg}{dooier} uitspreiden in plaats van een holte
te omsluiten.
\end{solution}

\begin{solution}{exo:b2:vertebrate-development:4}
\omterm{def:b2:vertebrate-development:induction}{Inductie} is de verandering van het lot van een competent weefsel door een
signaal uit een naburig weefsel. Vegetatieve cellen induceren mesoderm in
de bovenliggende marginale zone; de \omterm{def:b2:vertebrate-development:induction}{organisator} (mesoderm van de dorsale
lip) induceert de neurale plaat in het dorsale ectoderm; het oogblaasje
induceert de lens in het ectoderm van de kop.
\end{solution}

\begin{solution}{exo:b2:vertebrate-development:5}
$2^{12} = 4096$ cellen met straal $1/2^{4} = \qty{62.5}{\micro m}$;
oppervlak vermenigvuldigd met $2^{4} = 16$. Het extra oppervlak is het
membraan van de epithelen die de gastrulatie tot huid en darm zal plooien,
en het oppervlak waarlangs de \omterm{def:b2:vertebrate-development:blastula}{blastula} met haar omgeving uitwisselt.
\end{solution}

\begin{solution}{exo:b2:vertebrate-development:6}
$10^{-4}\times 2^{k} = 0.4$: $2^{k} = 4000$, $k = 12$. Met vier keer
zoveel DNA begint de verhouding bij $4\times 10^{-4}$ en bereikt ze $0.4$
bij $2^{k} = 1000$: de tiende deling, twee eerder.
\end{solution}

\begin{solution}{exo:b2:vertebrate-development:7}
$k = \ln 2/600 = \qty{1.16e-3}{s^{-1}}$; $\lambda = \sqrt{2\times
10^{-12}/1.16\times 10^{-3}} = \qty{42}{\micro m}$. Drempels bij
$x = \lambda\ln 2 = \qty{29}{\micro m}$ en $\lambda\ln 10 =
\qty{96}{\micro m}$: banden van 29, 67 en de rest.
\end{solution}

\begin{solution}{exo:b2:vertebrate-development:8}
Verdubbeling van de bron verschuift beide grenzen naar buiten met $\lambda\ln
2 = \qty{29}{\micro m}$: de eerste band verdubbelt, de tweede behoudt
haar breedte. Verdubbeling van $k$ deelt $\lambda$ door $\sqrt 2$
(\qty{29}{\micro m}) en, omdat $c_0 = j/\sqrt{Dk}$, deelt ook $c_0$ door
$\sqrt 2$: de grenzen gaan naar $\lambda'\ln(c_0'/T) =
29\times(0.69 - 0.35) = \qty{10}{\micro m}$ en $29\times(2.30 -
0.35) = \qty{57}{\micro m}$ --- beide banden krimpen.
\end{solution}

\begin{solution}{exo:b2:vertebrate-development:9}
Een dorsale lip uit een ongepigmenteerde gastrula van een watersalamander
werd naar de ventrale kant van een gepigmenteerde getransplanteerd; hij
rolde naar binnen en de gastheer vormde een tweede embryo --- \omterm{prop:b2:vertebrate-development:neurulation}{neurale buis},
somieten, darm --- vast aan het eerste. Het verschil in pigment liet zien
dat de tweede as uit cellen van de gastheer bestond, zodat het
transplantaat ze had geïnduceerd in plaats van de as zelf te bouwen. Het
sloot zelfdifferentiatie van het transplantaat als verklaring uit, en
toonde aan dat een klein gebied het signaal draagt dat de hele as
organiseert.
\end{solution}

\begin{solution}{exo:b2:vertebrate-development:10}
(a) Buikhuid: het vroege ectoderm is competent maar nog niet geïnduceerd,
en volgt dus zijn nieuwe omgeving. (b) Zenuwweefsel: in de late gastrula is
het geïnduceerd en gedetermineerd. (c) Neurale plaat: buikectoderm van een
late gastrula is nog competent, en de chorda induceert het. (d) Geen
dorsale structuren: een bol buikweefsel, omdat niets het mesoderm tot
dorsale lotbestemmingen of het ectoderm tot neurale induceert.
\end{solution}

\begin{solution}{exo:b2:vertebrate-development:11}
Bij \qty{18}{\celsius}: $90 + 11\times 30 = \qty{420}{min}$, 7 uur.
Bij \qty{10}{\celsius}: \qty{17.5}{h}. De transitie valt in beide gevallen
bij de twaalfde deling, omdat wat haar in gang zet de verhouding van DNA
tot cytoplasma is, die afhangt van het aantal delingen en niet van hoe lang
ze duurden: de klok telt verdubbelingen, geen uren.
\end{solution}

\begin{solution}{exo:b2:vertebrate-development:12}
Bij regulatieve ontwikkeling krijgen cellen hun lot door inductie vanuit
hun buren, zodat het embryo zich na verlies kan herstellen (tweelingen uit
een gesplitst ei, regeneratie van een getransplanteerd gebied); bij
mozaïekontwikkeling wordt het lot met plaatselijke cytoplasmatische
determinanten geërfd, zodat elke blastomeer alleen zijn eigen deel maakt en
een gesplitst embryo twee halve larven maakt. Het verschil is gradueel: de
grijze halvemaan van de kikker is een determinant, en mozaïekembryo's
gebruiken later inducties; elk embryo mengt erfenis en gesprek, en
gewervelden steunen langer op het gesprek.
\end{solution}

\begin{solution}{pb:b2:vertebrate-development:1}
\textbf{1.} Ei $\tfrac43\pi(0.7)^{3} = \qty{1.44}{mm^3}$; kern
$\tfrac43\pi(0.03)^{3} = \qty{1.1e-4}{mm^3}$; verhouding $8\times
10^{-5}$.
\textbf{2.} $2^{k} = 4000$: $k = 12$, 4096 cellen.
\textbf{3.} $75 + 11\times 30 = \qty{405}{min}$, ongeveer 6.75 uur.
\textbf{4.} Straal $0.7/2^{4} = \qty{44}{\micro m}$; verhouding $8\times
10^{-5}\times 4096 = 0.32$ --- die van een gewone cel: de kernen hebben het
cytoplasma ingehaald.
\textbf{5.} $2^{4} = 16$.
\textbf{6.} Ongeveer $4095\,c$ aan DNA; het ei bevat genoeg opgeslagen
histonen, polymerasen en nucleotiden voor twaalf delingen, zodat er geen
transcriptie nodig is.
\textbf{7.} Zygotische genen schakelen aan, de delingen vertragen en
verliezen hun synchronie, cellen worden beweeglijk. Het injecteren van extra
DNA vervroegt de transitie met zoveel delingen als het DNA in overmaat is,
wat een teller van delingen of van tijd niet zou kunnen verklaren.
\textbf{8.} Tegenover het binnenkomstpunt van de zaadcel, bij de grijze
halvemaan, vastgelegd door de rotatie van de eicortex na de bevruchting.
\textbf{9.} Eerst dorsaal mesoderm (prechordale plaat, dan chorda) op het
dak van de oerdarm in de middellijn; lateraal mesoderm (somieten, laterale
plaat) ernaast; het endoderm bekleedt de oerdarm; het ectoderm spreidt zich
door epibolie over de buitenkant uit.
\textbf{10.} $\qty{2000}{\micro m}/\qty{600}{min} = \qty{3.3}{\micro
m/min}$: drie keer een fibroblast --- een gecoördineerde laag beweegt
sneller dan een losse cel.
\textbf{11.} De celbewegingen hebben zygotische genproducten nodig (nieuwe
adhesiemoleculen, beweeglijkheid) en vertraagde cycli; vóór de transitie
kunnen de cellen alleen delen.
\textbf{12.} Elk chordadier heeft in een of ander stadium een chorda, en
gewervelden bouwen de wervelkolom eromheen; de \omterm{prop:b2:vertebrate-development:neurulation}{neurale buis} wordt erdoor
geïnduceerd en is het gevolg, niet de oorsprong, van het plan.
\textbf{13.} Zonder samentrekking van de flescellen vormt zich geen
blastoporus, rolt er niets naar binnen, en blijft het embryo een bol met
zijn lagen op hun plaats: geen darm, geen chorda, geen as.
\textbf{14.} $k = \ln 2/1200 = \qty{5.8e-4}{s^{-1}}$; $\lambda =
\sqrt{10^{-12}/5.8\times 10^{-4}} = \qty{42}{\micro m}$.
\textbf{15.} $x_1 = 42\ln 2.5 = \qty{38}{\micro m}$; $x_2 = 42\ln 12.5
= \qty{105}{\micro m}$.
\textbf{16.} Banden van \qty{38}{\micro m} (2.5 cellen), \qty{67}{\micro
m} (4.5 cellen), en de rest.
\textbf{17.} $x_1 = 42\ln 1.25 = \qty{9}{\micro m}$; $x_2 = 42\ln 6.25
= \qty{77}{\micro m}$: beide grenzen schuiven naar binnen op met $\lambda\ln 2 =
\qty{29}{\micro m}$; de eerste band krimpt van 38 tot 9, de middelste band
behoudt zijn breedte, de buitenste band groeit.
\textbf{18.} Verplaatst vóór het aflezen wordt hij lot A --- hij leest de
concentratie waar hij nu is; verplaatst na het aflezen behoudt hij lot B
--- hij was gedetermineerd.
\textbf{19.} $x_T = \lambda\ln(c_0/T)$: een verandering van de bron met een
factor $f$ verschuift elke grens met $\lambda\ln f$, zodat een fout van een
factor twee het patroon maar $0.7\lambda$ verschuift, een fractie van een
veld dat meerdere $\lambda$ beslaat --- het patroon is vrijwel ongevoelig
voor de precieze hoeveelheid morfogeen.
\textbf{20.} Dorsale lip naar ventraal: een tweede as. Ventrale marginale
zone naar dorsaal: ze wordt door haar omgeving opnieuw gespecificeerd en
maakt dorsaal weefsel; niets extra.
\textbf{21.} In een neurula is het ectoderm niet meer competent: het
transplantaat maakt chorda maar induceert geen tweede \omterm{prop:b2:vertebrate-development:neurulation}{neurale buis}.
\textbf{22.} Gesplitst in het vlak van de halvemaan: twee volledige
embryo's; loodrecht daarop: één embryo en een bol buikweefsel.
\textbf{23.} Alleen: een bol opperhuid. Met vegetatieve cellen: mesoderm
--- spier, chorda --- geïnduceerd in het kapje.
\textbf{24.} Om te bewijzen dat de tweede as in cellen van de gastheer was
geïnduceerd en niet door het transplantaat was gebouwd; was het
transplantaat niet te onderscheiden geweest, dan had zelfdifferentiatie
niet kunnen worden uitgesloten. Een transplantaat van een andere soort
werkt (ze gebruikten twee soorten watersalamanders), en dat is precies wat
de cellen te onderscheiden maakte.
\textbf{25.} Twaalfde \omterm{def:b2:vertebrate-development:egg}{klieving}, 4096 cellen, na ongeveer 6.75 uur; banden
van 38 en \qty{67}{\micro m}, en alles voorbij \qty{105}{\micro
m}.
\end{solution}
